\chapter{RESULTS AND DISCUSSION}

\section{INTRODUCTION}
This chapter presents the outcomes achieved following the implementation and testing of the ScrapKart AI platform. It evaluates the performance of the AI classification model, analyzes the efficiency gains in the scrap collection process, and discusses the overall user experience based on the User Acceptance Testing (UAT) phase.

\section{AI MODEL PERFORMANCE}
The core technological achievement of this project is the AI-driven scrap classification engine based on the YOLOv8 architecture. To evaluate its effectiveness, the model was tested against a separate validation dataset consisting of 2,000 diverse images not seen during the training phase.

\begin{itemize}
    \item \textbf{Accuracy Metrics:} The model achieved an overall \textbf{Mean Average Precision (mAP) of 89.4\%} at an Intersection over Union (IoU) threshold of 0.5. 
    \item \textbf{Category Specific Performance:} The model demonstrated the highest accuracy in identifying "Cardboard/Paper" (94\% precision) and "Plastic Bottles" (91\% precision). "Electronic Waste (E-waste)" showed a slightly lower precision of 82\%, primarily due to the high visual variance and complexity of electronic components.
    \item \textbf{Inference Speed:} Operating on the cloud GPU infrastructure, the average inference time per image was recorded at 120 milliseconds. This rapid processing ensures a near-instantaneous user experience, where the customer receives a price quotation almost immediately after uploading a photo.
\end{itemize}

\begin{figure}[H]
    \centering
    % \includegraphics[width=0.8\textwidth]{Figures/Confusion_Matrix.png}
    \vspace{5cm}
    \caption{Confusion Matrix of AI Scrap Classification}
    \label{fig:confusion_matrix}
\end{figure}

\section{OPERATIONAL EFFICIENCY AND ROUTE OPTIMIZATION}
A primary objective was to improve the logistics of scrap collectors. During the UAT phase, telemetry data from 5 participating collectors was analyzed over a two-week period.
\begin{itemize}
    \item \textbf{Distance Reduction:} By utilizing the Google Maps integrated route optimization algorithm, collectors experienced an average \textbf{22\% reduction in total daily travel distance} compared to their traditional, unplanned collection routes.
    \item \textbf{Time Efficiency:} The time spent actively searching for scrap was eliminated. Collectors only traveled to confirmed locations, resulting in a \textbf{35\% increase in the volume of scrap collected per hour}.
\end{itemize}

\section{USER ACCEPTANCE AND INTERFACE EVALUATION}
Feedback collected from the 20 test households indicated high satisfaction with the transparency of the platform.
\begin{itemize}
    \item \textbf{Transparency:} 90\% of users reported that the AI-generated price quotation provided them with confidence that they were receiving a fair market value for their waste, a significant improvement over traditional haggling.
    \item \textbf{Sustainability Impact:} Users responded highly positively to the Sustainability Dashboard. Displaying metrics such as "Equivalent Trees Saved" and "Carbon Emissions Prevented" actively motivated 60\% of the test group to upload additional scrap items within the two-week testing period.
\end{itemize}

\begin{figure}[H]
    \centering
    % \includegraphics[width=0.8\textwidth]{Figures/Sustainability_Dashboard.png}
    \vspace{5cm}
    \caption{User Interface: Sustainability Dashboard Output}
    \label{fig:ui_dashboard}
\end{figure}

\section{ADVANTAGES ACHIEVED}
The practical implementation of ScrapKart AI successfully realized the theoretical advantages outlined in Chapter 1:
\begin{enumerate}
    \item \textbf{Digitization of Informal Sector:} Successfully transitioned an unorganized process into a structured, data-driven digital ecosystem.
    \item \textbf{Fair Economics:} Automated the valuation process, eliminating human bias and ensuring fair compensation based on real-time data.
    \item \textbf{Environmental Optimization:} Reduced the carbon footprint of collectors through smart routing and increased overall recycling engagement among households.
\end{enumerate}

\section{LIMITATIONS}
Despite the successes, certain limitations were observed during deployment:
\begin{itemize}
    \item \textbf{Weight Estimation Accuracy:} While the AI excels at classifying materials, estimating the exact weight (and therefore the exact price) from a 2D image remains challenging. Hidden objects or dense materials can lead to slight discrepancies between the AI quotation and the final physical measurement.
    \item \textbf{Internet Dependency:} Both customers and collectors require a stable mobile internet connection to access the platform, upload high-quality images, and utilize the live navigation features, which may be a constraint in areas with poor network coverage.
\end{itemize}
